\documentclass[11pt]{article}
\usepackage[T1]{fontenc}
\usepackage[utf8]{inputenc}
\usepackage[english]{babel}
\usepackage{amsmath,amssymb,amsthm,booktabs}
\usepackage[margin=1in]{geometry}
\usepackage[colorlinks=true,linkcolor=black,citecolor=blue,urlcolor=blue]{hyperref}
\usepackage{url}
\hypersetup{pdftitle={Extending A358321 by prime-power enumeration and arithmetic exclusions},pdfauthor={Carlo Corti}}
\newtheorem{theorem}{Theorem}
\newtheorem{proposition}[theorem]{Proposition}
\newtheorem{lemma}[theorem]{Lemma}
\newtheorem{corollary}[theorem]{Corollary}
\theoremstyle{remark}
\newtheorem{remark}[theorem]{Remark}
\title{Extending A358321 by prime-power enumeration\protect\\and arithmetic exclusions}
\author{Carlo Corti\\Independent researcher\\\texttt{carlo.corti@outlook.com}}
\date{November 2026}
\begin{document}
\maketitle
\begin{abstract}
For $n\geq3$, let $a(n)$ be the least positive index $k$ for which the
$n$-gonal number $P_n(k)=k((n-2)k-(n-4))/2$ has exactly $n$ distinct prime
factors. Subject to the recorded producer coverage and exact arithmetic exclusions
described below, we report four locally validated terms of OEIS A358321
for $n=16,17,18,19$, with
$a(16)=119535945060$, $a(17)=1113316946280$,
$a(18)=4278672701330$, and $a(19)=153290693815740$.
The search constructs products of increasing distinct prime bases with
arbitrary positive exponents, following the constructive method published
by Daniel Suteu. Minimality rests on the recorded exhaustive producer
coverage; for $n=19$ it also uses exact exclusions of unfinished subtrees.
Indeed, for $p=3$ or $5$, divisibility of $P_{19}(k)$ by $p$ forces
divisibility by $p^2$. Applied after the search, this observation completes
the mathematical coverage without asserting native program completion.
Independent Python and PARI/GP certificates verify the positive witnesses,
not a second exhaustive negative enumeration. A separate final interval
$(166539042068893,170000000000000]$ contains no matching index under
the same execution and coverage assumptions.
The corresponding four polygonal values give same-index consequences for
A358862; neither this interval nor those consequences determine $a(20)$.
\end{abstract}

\noindent\textbf{2020 Mathematics Subject Classification.}
Primary 11Y55; Secondary 11A51, 11B83.

\medskip
\noindent\textbf{Key words and phrases.}
polygonal number, distinct prime factors, integer sequence, exhaustive enumeration,
computational minimality, OEIS A358321.

\section{Definition, provenance and contribution}\label{sec:intro}
For integers $n\geq3$ and $k\geq1$, write
\begin{equation}\label{eq:polygon}
 P_n(k)=\frac{k((n-2)k-(n-4))}{2}
       =k+(n-2)\frac{k(k-1)}2.
\end{equation}
This is the standard $n$-gonal number formula~\cite{polygonal}.
For a positive integer $v$, let $\omega(v)$ count its distinct prime divisors,
with $\omega(1)=0$; $\Omega(v)$ counts them with multiplicity~\cite{distinct}.
The target sequence is
\begin{equation}\label{eq:def}
 a(n)=\min\{k\geq1:\omega(P_n(k))=n\},
\end{equation}
whenever the set is nonempty. Thus the parameter $n$ controls both the
polygonal order and the required support cardinality. For example,
$P_3(11)=66=2\cdot3\cdot11$ and $a(3)=11$.

A358321 was submitted by Ilya Gutkovskiy on 12 December 2022~\cite{oeis321}.
The source prefix consulted before the computation comprised $n=3,\ldots,15$:
\begin{multline}\label{eq:prefix}
11,210,87,228,1155,7854,66612,395646,2193303,8389010,\\
122574155,630341910,6066475415.
\end{multline}
The official OEIS export retrieved through the OEIS-maintained GitHub
repository on 9 November 2026 exposes the same prefix; a check on
10 November returned identical entry text. Direct retrieval of the
corresponding live b-files and detailed editorial history was unavailable
in the author's retrieval environment. The contribution asserted here is a
local validated extension of that stated baseline, without a claim of
worldwide priority or an assertion that OEIS has incorporated these terms.

The closely related sequence A358862 records the least \emph{value}, rather
than its index~\cite{oeis862}. Its entry credits Gutkovskiy, 3 December 2022,
and Daniel Suteu for terms $12$--$15$ and two PARI programs dated
4 December 2022. Suteu's constructive program enumerates integers with a
specified number of distinct prime factors by increasing prime bases and
positive powers, then tests polygonality. Its driver also searches successive adjacent value windows. This
constructive enumeration and windowing are the central algorithmic
prior art used here. The same entry contains Harvey P. Dale's Mathematica
program of 9 September 2023 for the first eight terms. The present work is
a fresh C implementation of that method, with exact tail bounds, streaming
minima and retained coverage evidence; it does not claim to have originated
constructive prime-power enumeration.

The distinctive feature of the last minimum is its mixed proof: the
computationally completed subtrees are combined with an elementary
valuation obstruction for the unfinished part. The obstruction was applied
after the expensive search; it was not an early subtree filter of the
producer. We first describe the finite search, then prove this obstruction
and state the computational result with its precise evidence boundary.

\section{A finite and nonredundant search}\label{sec:method}
\subsection{Bounds and exact inversion}
Let $q_i$ be the $i$th prime and put $Q_r=\prod_{i=1}^r q_i$, with $Q_0=1$.
If $\omega(v)=n$, then $v\geq Q_n$. Moreover every prime divisor $p$ of such
a $v\leq B$ satisfies
\begin{equation}\label{eq:cap}
 p\leq\left\lfloor B/Q_{n-1}\right\rfloor,
\end{equation}
because the other $n-1$ distinct primes have product at least $Q_{n-1}$.
The bound remains valid when some exponents exceed one.
For an index cutoff $K$, it is therefore sufficient to construct the
$n$-support integers in $[Q_n,P_n(K)]$.

For fixed $n$, direct subtraction gives
\begin{equation}\label{eq:monotone}
 P_n(k+1)-P_n(k)=(n-2)k+1>0.
\end{equation}
A positive integer $v$ is $n$-gonal exactly when
\begin{equation}\label{eq:inverse}
 D=(n-4)^2+8(n-2)v=s^2,\qquad
 k=\frac{s+n-4}{2(n-2)}\in\mathbb Z_{\geq1},
\end{equation}
where $s$ is the nonnegative integer square root. This follows by completing
the square in~\eqref{eq:polygon}; the same inverse criterion is given
in~\cite{polygonal}. A square discriminant alone is insufficient: the
divisibility and positive-index tests are required. The implementation also
substitutes the recovered $k$ into~\eqref{eq:polygon} before accepting it.
For $n=16$, the equivalent reduced formulas
\begin{equation}\label{eq:16}
 P_{16}(k)=k(7k-6),\qquad t^2=7v+9,\qquad k=(t+3)/7
\end{equation}
reduce the size of the radicand. Here $t=s/4$ when $s$ is the
nonnegative root in~\eqref{eq:inverse} for $n=16$.

\subsection{Increasing prime powers and tail pruning}
Every integer with $n$ distinct prime divisors has a unique representation
\begin{equation}\label{eq:tuple}
 v=p_1^{e_1}\cdots p_n^{e_n},\qquad
 p_1<\cdots<p_n,\qquad e_i\geq1.
\end{equation}
At a recursive node, $m$ is the current product and $r$ prime bases remain
to be selected. For each next prime $p$ in increasing order, the generator
tries its powers with exponent $1,2,\ldots$, and recurses using only primes
larger than $p$. If $t_1<\cdots<t_r$ are the least allowed remaining primes,
then every completion is at least
\begin{equation}\label{eq:tail}
 m\prod_{i=1}^r t_i.
\end{equation}
A branch is discarded when this lower bound exceeds $B$. The bound is
monotone as the next prime or its exponent increases, justifying termination
of the corresponding loop. The prime list is obtained by a sieve through
the inclusive cap~\eqref{eq:cap}. Products below a window's lower bound
may be constructed, but are not tested for membership in that window.

\begin{proposition}[Completeness of the finite generator]\label{prop:generator}
For a finite positive bound $B$, the preceding procedure generates each
integer $v\leq B$ with $\omega(v)=n$ exactly once, before the window lower
bound and polygonality tests.
\end{proposition}
\begin{proof}
The prime cap contains every prime of every admissible tuple by~\eqref{eq:cap}.
Fix one such tuple. At any of its prefixes, the least-tail product in
\eqref{eq:tail} cannot exceed the actual completion, so its path is never
pruned. Monotonicity of the tail bound ensures that every earlier prime or
exponent choice also passes its relevant bound, so no preceding loop break
prevents reaching the prescribed choice at any depth. Induction on the remaining depth reaches its leaf.
Conversely, every leaf has exactly $n$ increasing prime bases with positive
exponents, so has support size $n$. Unique factorization makes two distinct
tuples unequal as integers. There are finitely many primes under the cap;
each exponent is bounded by $p^e\leq B$; recursion has depth $n$. Thus all
loops terminate. The empty search case $B<Q_n$ follows from the primorial
bound and requires no recursive leaves.
\end{proof}

Separate value windows have inclusive lower bound one greater than the
previous upper bound. For a fixed $n$, their coverage therefore joins without
a gap or overlap. Since tuple order is not numerical product order, the
producer retains a minimum over the whole completed window rather than
taking the first polygonal leaf. Equalities~\eqref{eq:monotone}
and~\eqref{eq:inverse} make least value and least index equivalent.

\subsection{Arithmetic and persistent minima}\label{sec:implementation}
The source uses C17 with GCC's \texttt{unsigned \_\_int128}
extension~\cite{gcc} for products, discriminants and wide counters; the
implementation stores prime bases as 32-bit unsigned integers. Before multiplying by $p$, a bound of the form
$m\leq B/p$ prevents overflow. The numerical admission checks bound both
the sieve cap by $2^{32}-1$ and the generic discriminant by $2^{128}-1$.
The implemented input domain $13\leq n\leq24$ is a numerical capacity
statement, not a computed range or a guarantee of feasible running time.

Necessary square-residue tests modulo $64$ and $63$ precede the exact root
calculation at a leaf. The masks are constructed from all square residues,
so rejecting a nonsquare residue cannot discard a polygonal value. A
floating-point root estimate affects speed only: exact integer-square
comparisons correct it in both directions and the acceptance test is integer
arithmetic. In particular, no probable-prime test is used to decide membership
of an enumerated leaf; its prime support is already explicit.

The tree is partitioned by fixed initial prime-power tuples. The final
$n=19$ jobs use three-base prefixes. OpenMP workers process these disjoint
units~\cite{gcc}; a unit may be interrupted and resumed from its last fully
processed tuple. In the later producer the cursor order is the lexical
order of prime/exponent choices used by the recursion, rather than the order
of products. An improved global minimum is durably recorded before a cursor
or completion record advances. An interruption before that advance may
repeat work, but cannot silently skip an unrecorded smaller witness. A
completion record certifies exhaustion of a unit's prescribed traversal;
a partial record does not.

The earlier producer, denoted E in Table~\ref{tab:windows}, executed all
six $n=16,17,18$ windows and the first three $n=19$ windows. It used the
same increasing-prime-power generation, depth-three partition, tail bounds,
leaf masks and exact membership criterion. A completed unit's minimum and
certificate were persisted together in its completion record; partial-unit
work was discarded and that unit restarted on the next call. The second
and third $n=19$ windows required such continuation. The later producer L
retains the tuple cursors and improving minimum records described above.
Its first $n=19$ window has initial index cutoff $81587135064138$ and starts
at value $P_{19}(81587135064138)+1$, exactly after E's final covered value.
Thus E accounts for nine rows and L for three rows of the table.
Both producers are retained with their own native verifiers; a separate
artifact auditor checks the corresponding histories and their exact join.
Integrity of these histories supports provenance, but the mathematical
coverage argument also requires the generator, the exact window boundaries,
the disjoint unit definitions and their completion or analytic exclusion.
The reproducibility material is distributed through the project repository
and Zenodo release history~\cite{corti}.

\section{Arithmetic exclusion of entire subtrees}\label{sec:arithmetic}
Put $L_n(k)=(n-2)k-(n-4)$, which is positive for $n\geq3$, $k\geq1$.
The two linear factors satisfy
\begin{equation}\label{eq:gcd}
 \gcd(k,L_n(k))=\gcd(k,n-4).
\end{equation}
Here $\gcd(k,-1)=1$ and $\gcd(k,0)=k$, so the identity includes $n=3,4$.
It follows by subtracting $(n-2)k$ inside the gcd. The shared prime
support is thus confined to the primes of $n-4$, except that $n=4$ is the
square case. For example~\eqref{eq:16} gives
\begin{equation}\label{eq:support16}
 \omega(P_{16}(k))=\omega(k)+\omega(7k-6)
                  -\boldsymbol{1}_{2\mid k}-\boldsymbol{1}_{3\mid k},
\end{equation}
because $\gcd(k,7k-6)=\gcd(k,6)$. The symbol
$\boldsymbol{1}_{E}$ is one when $E$ holds and zero otherwise.

For a prime $p$ and a positive integer $x$, write $v_p(x)$ for the exponent
of $p$ in $x$.
\begin{lemma}[Odd-prime obstruction]\label{lem:valuation}
Let $n\geq3$ and let $p$ be an odd prime dividing $n-4$. For every $k\geq1$,
\begin{equation}\label{eq:noone}
 p\mid P_n(k)\quad\Longleftrightarrow\quad p\mid k,
 \qquad p\mid P_n(k)\quad\Longrightarrow\quad p^2\mid P_n(k).
\end{equation}
In particular $v_p(P_n(k))\ne1$.
\end{lemma}
\begin{proof}
Since $n-2\equiv2\pmod p$ and $2$ is invertible modulo $p$,
$P_n(k)\equiv k^2\pmod p$. Thus $p\mid P_n(k)$ if and only if $p\mid k$.
In that case both $k$ and $L_n(k)$ are divisible by $p$, and division of
their product by $2$ does not remove a factor of $p$. This proves the claim.
\end{proof}

The restriction to odd primes is essential to this proof. For example,
$P_{10}(2)=10$ has 2-adic valuation one although $2\mid10-4$.
For $n=3$ there is no prime divisor of $n-4=-1$. For $n=4$ every odd prime
divides $n-4=0$ and the conclusion agrees with $P_4(k)=k^2$; in that case
all prime valuations, including the 2-adic valuation, are even. These are
elementary consequences of the defining polynomial; no novelty claim is
made for the gcd identity or the lemma.

\subsection{The two relevant primes for \texorpdfstring{$n=19$}{n=19}}
Here
\begin{equation}\label{eq:19}
 P_{19}(k)=k(17k-15)/2,\qquad n-4=15.
\end{equation}
Lemma~\ref{lem:valuation} applies to $p=3,5$. As finite arithmetic checks,
the residue lists for $k=0,\ldots,8$ modulo $9$, and for
$k=0,\ldots,24$ modulo $25$, are respectively
\begin{align}
 &(0,1,1,0,7,4,0,4,7),\label{eq:mod9}\\
 &\begin{aligned}[t]
 (0,1,19,4,6,0,11,14,9,21,0,21,9,\qquad\\
 14,11,0,6,4,19,1,0,16,24,24,16).
 \end{aligned}\label{eq:mod25}
\end{align}
The second display is one list continued on its next line.
Neither $3,6$ modulo $9$ nor $5,10,15,20$ modulo $25$ occurs. These
lists apply to every integer $k$ by the exact identities
\begin{align}
 P_{19}(k+9)-P_{19}(k)&=9(17k+69),\label{eq:period9}\\
 P_{19}(k+25)-P_{19}(k)&=25(17k+205).\label{eq:period25}
\end{align}
The direct divisibility proof establishes the universal assertion without
relying on numerical sampling. Python and PARI/GP also check the displayed
lists and period identities independently.

If a fixed prefix contains $3^1$ or $5^1$, every descendant has that exact
valuation: later prime bases are strictly greater than those already fixed.
Its entire subtree is therefore empty of 19-gonal values, regardless of how
much of the subtree the native program processed. This argument depends on
the prime being fixed in the prefix; it does not reject a branch simply
because that prime might occur later.

Unit IDs are local to each window: zero-based numbering follows the
lexicographic order of $(p_1,e_1,p_2,e_2,p_3,e_3)$ in its prefix map.
Changing the upper bound can insert prefixes and change later IDs.
In the decisive window the sole unfinished unit, numbered $0$,
fixes $(2^1,3^1,5^1)$.
In the final supplement the unfinished units $0$ and $110$ fix
$(2^1,3^1,5^1)$ and $(2^2,3^1,5^1)$, respectively. All three subtrees are
excluded by the lemma. The exponent mapping is checked against the actual
prefix definitions and stored cursors, not inferred from a unit's number.
The $p=3$ case alone suffices to exclude all three pending subtrees;
$p=5$ gives an additional sufficient obstruction for them. Across the
full prefix maps, the union of the two obstruction classes contains
$276$ of $577$ and $279$ of $590$ units, respectively. Most of these units
were already completed computationally: only the pending units are added
as analytic coverage in the proof, so completed units are never counted twice.

The $p=3$ exclusion was adopted after the search had produced the $n=19$
witness and had left one native unit unfinished. The $p=5$ instance was
then documented before the final supplement. The unchanged producer still
applied its residue masks at leaves; it did not implement these analytic
prefix exclusions. Mathematical closure and native completion therefore
remain distinct: the last two jobs are retained as incomplete and unsealed.

\section{Coverage, certificates and results}\label{sec:results}
\subsection{Finite coverage and computational theorem}
Table~\ref{tab:windows} gives the inclusive index cutoffs and prefix counts.
For each $n$, the first value window starts at $Q_n$, and each later window
starts at the previous $P_n(K)+1$. For $n=19$ this includes the transition
between the earlier and later producers. No index $k$ with $P_n(k)<Q_n$ can satisfy the predicate, since the
primorial threshold bounds polygonal values.

\begin{table}[htbp]
\centering
\small
\begin{tabular}{crrrrl}
\toprule
Producer & $n$ & $K$ & Closed & Total & Minimum / analytic addition\\
\midrule
E & 16 & 68231913854 & 158 & 158 & none\\
E & 16 & 152571197612 & 320 & 320 & 119535945060\\
E & 17 & 506328003064 & 154 & 154 & none\\
E & 17 & 876985826602 & 251 & 251 & none\\
E & 17 & 1518984009193 & 393 & 393 & 1113316946280\\
E & 18 & 9762008494901 & 337 & 337 & 4278672701330\\
\midrule
E & 19 & 67989279220115 & 297 & 297 & none\\
E & 19 & 77519705336978 & 332 & 332 & none\\
E & 19 & 81587135064138 & 343 & 343 & none\\
L & 19 & 117760885979226 & 451 & 451 & none\\
L & 19 & 166539042068893 & 576 & 577 & 153290693815740; unit $0$ empty\\
L & 19 & 170000000000000 & 588 & 590 & none; units $0,110$ empty\\
\bottomrule
\end{tabular}
\caption{Retained coverage. E and L denote the earlier and later producer.
Closed counts distinct natively completed prefix units within each window. The final two rows have full mathematical coverage only after
adding the stated analytic exclusions; their native jobs remain incomplete.}
\label{tab:windows}
\end{table}

\begin{theorem}[Validated computational minima]\label{thm:main}
Subject to the correct execution of the validated prime-power producers
and the recorded disjoint coverage in Table~\ref{tab:windows}, together with
the exact prefix exclusions of Lemma~\ref{lem:valuation}, the minima in
\eqref{eq:def} for $n=16,17,18,19$ and their polygonal values are those in
Table~\ref{tab:values}. For $n=19$, no $k$ in the separate interval
\begin{equation}\label{eq:supplement}
 166539042068893<k\leq170000000000000
\end{equation}
satisfies $\omega(P_{19}(k))=19$.
\end{theorem}
\begin{proof}[Computational proof with arithmetic exclusions]
Under the execution and coverage hypotheses of the theorem, for $n=16$
all $158+320=478$ units are complete and their minimum is the
reported index. For $n=17$ the corresponding counts are
$154+251+393=798$, and for $n=18$ all $337$ units are complete. The
primorial bound, the adjacent value windows and
Proposition~\ref{prop:generator} account for every possible smaller value.
Independent positive certificates establish membership at the reported
indices.

For $n=19$, the first four rows for that $n$ have no hits and complete
native coverage through $K=117760885979226$. The next window has $576$
natively complete units out of $577$, and the durably recorded global
minimum over the union of their processed tuples is $153290693815740$. Its positive
certificate belongs to complete unit $32$, whose three-base prefix is
$(2^1,3^3,5^2)$. The only remaining unit is $0$, excluded by
Lemma~\ref{lem:valuation}. Thus no smaller matching value remains outside
the proved coverage. Strict monotonicity~\eqref{eq:monotone} establishes
the least index.

The supplement starts at value $P_{19}(166539042068893)+1$ and has $588$
complete native units out of $590$, with no positive minimum. Its only
pending units $0$ and $110$ are also excluded by the lemma. This proves
\eqref{eq:supplement} empty of matching indices. The native records remain
$576/577$ and $588/590$; no completion record is supplied in place of the
arithmetic proof. The negative computational part is not a second,
independently implemented exhaustive enumeration.
\end{proof}

\begin{table}[htbp]
\centering
\small
\begin{tabular}{rrr}
\toprule
$n$ & $a(n)$ & $P_n(a(n))$\\
\midrule
16 & 119535945060 & 100021895128994153154840\\
17 & 1113316946280 & 9296059671549446483637180\\
18 & 4278672701330 & 146456320680822524361241890\\
19 & 153290693815740 & 199733312889341920908906236550\\
\bottomrule
\end{tabular}
\caption{The four locally admitted terms and the exact same-index polygonal
values. The last column is derived from the middle column, rather than
from a separate enumeration.}\label{tab:values}
\end{table}

The supplement records an additional finite no-hit interval after the minimum
was established; it is not needed to prove that minimum.
The producer records minima, rather than all matching indices. In particular,
the proof of the least $n=19$ index does not enumerate every matching index
between that minimum and the supplement's lower cutoff. Hence the empty
supplement must not be enlarged to the interval $(a(19),170000000000000]$
without additional evidence. The cutoff concerns $P_{19}$, and gives no
search coverage for $P_{20}$.

\subsection{Independent positive certificates and their limits}\label{sec:cert}
For each reported term, an independent Python checker factors the two linear
factors by elementary trial division, combines their exponents and accounts
for division by $2$ in~\eqref{eq:polygon}. For even orders it can instead
factor the already simplified product, as in~\eqref{eq:16}. It compares the
prime support and full product with the claimed value, and checks the exact
inverse discriminant. Changed-index, changed-value and changed-factor
controls are rejected. Separate PARI/GP inputs factor the values and verify
primality, support size and inversion. The relevant meanings of
\texttt{omega}, \texttt{bigomega}, \texttt{factor}, \texttt{isprime} and
\texttt{ispolygonal} are documented in the official manual~\cite{pari}.
Explicit primality verification avoids treating an unproved large factor
from a general factorization routine as automatically prime.

For $k=153290693815740$, the retained certificate is
\begin{equation}\label{eq:factor19}
\begin{aligned}
 P_{19}(k)={}&2\cdot3^3\cdot5^2\cdot7\cdot11\cdot13\cdot17^2\cdot19
                 \cdot23\cdot29\cdot31\cdot37\\
             &{}\cdot41\cdot43\cdot67\cdot71\cdot73\cdot101\cdot569.
\end{aligned}
\end{equation}
There are $19$ distinct bases. In particular the valuations at $3$ and $5$
are $3$ and $2$, consistent with the arithmetic obstruction. The Python
check also factors the two linear factors independently of the producer's
prime-power traversal. PARI/GP provides another arithmetic backend. Both
certify this positive witness, not the assertion that every smaller index
fails the predicate.

A retained Python test harness constructs a diagnostic C fixture from the
producer and compares its complete small output sets for orders $n=4,5,6,7$
against a direct distinct-factor-count sieve, with filters both enabled and
disabled. These diagnostic orders lie outside the production input domain. This
is distinct from the built-in self-test. Further retained tests cover
resumed tails, malformed records and interruption/recovery. The small-set oracle counts
prime divisors by marking multiples, rather than by the producer's
prime-power recursion. These finite tests do not independently enumerate
the full production ranges. Replaying Suteu's published GP method on the
13 source terms checks its implementation and the value/index convention;
it is reproduction of the same constructive method, not independent
mathematical evidence for all new non-hits. Similarly, an artifact auditor
checks integrity, counters, prefix identities, positive products and
coverage joins; it does not itself regenerate every omitted leaf.

The resulting certification boundary has three components: mathematical
completeness of the finite generator, correct execution and coverage of its
computationally closed subtrees, and independently checked arithmetic of
the positive witnesses and excluded prefixes. The conditional execution
basis in Theorem~\ref{thm:main} makes the residual common-mode risk explicit.

\section{Consequences for related sequences}\label{sec:related}
Write $b(n)$ for A358862. Since~\eqref{eq:monotone} is strictly positive,
\begin{equation}\label{eq:relation}
 b(n)=P_n(a(n))
\end{equation}
whenever the minimum exists. Conversely, the exact inverse~\eqref{eq:inverse}
recovers $a(n)$ from $b(n)$. All 13 source pairs for $n=3,\ldots,15$
satisfy this relation; for example, $P_3(11)=66$ and
$P_{15}(6066475415)=239213805711830629680$. Thus the final column of
Table~\ref{tab:values} is also the locally validated same-index extension
of A358862 through $19$. It requires no second autonomous search.

A359014, also submitted by Gutkovskiy on 12 December 2022, instead minimizes
the index under $\Omega(P_n(k))=n$; its entry includes Suteu's PARI program
and credits his terms $29$--$34$~\cite{oeis014}. This distinction changes
the minima already at $n=3$: $P_3(7)=28=2^2\cdot7$ has $\Omega=3$ and
$\omega=2$. Its longer source table therefore cannot extend the present
sequence. More generally, varying $n$ changes both the polynomial and the
target support, so none of the relations here converts the finite $n=19$
coverage into a determination of $a(20)$. No assertion of existence for
every $n$ is implied by the finite results.

\section{Implementation context and reproducibility}\label{sec:repro}
The computation was run on an AMD Ryzen 9 7940HS workstation with eight
physical cores and sixteen hardware threads, under Linux Mint 22.3
(x86-64; kernel 6.17.0-42-generic). The production build used GCC 13.3.0,
C17, OpenMP and the flags
\begin{quote}\small\ttfamily\raggedright
-O3 -march=native -mtune=native -flto -funroll-loops
-fomit-frame-pointer -std=c17 -fopenmp
\end{quote}
with the retained warning options and the mathematics library. These flags
and the original build specification are retained; \texttt{-march=native}
means that the executable itself is specific to the build host. The final
jobs were configured with eight worker threads. This is a configuration,
not an eightfold speedup claim. The checkers use Python 3.12.3 standard-library
integer arithmetic and PARI/GP 2.15.4. Neither Sidef nor a symbolic algebra
package is a dependency of the production enumeration.

As one scoped measurement, the final interval~\eqref{eq:supplement} consumed
about $9900$ seconds of accumulated native search wall time in three calls,
including work on prefixes covered analytically but not skipped by
the unchanged producer. This is not
CPU time or a conservative resource reservation; it also excludes elapsed
time between calls. Its stored leaf count is $2791504259667$, and its
stored exact-root-call count is $52068962797$. The lower value cutoff skips
membership tests but does not eliminate construction of below-cutoff tuples,
so these counters are not counts of distinct admissible indices in the
supplement. They must not be used to extrapolate a guaranteed runtime for
another order or horizon. No single whole-project timing or peak-memory
benchmark is asserted here.

The local reproducibility material retains the source and build specification,
the earlier producer needed for its own histories, the 17-term table, exact
window and prefix records, positive Python/PARI certificates, modular
certificates and validation programs with their inputs and execution outputs.
A reconstruction of the proof checks, in order, the historical prefix and
value/index map, the producer identities appropriate to each history, the
coverage joins and completed units, the analytically excluded prefixes, and
the positive factorizations. These steps check the retained computation;
a full independent reproduction of the negative search is a stronger and
separate task. The repository and archive reference~\cite{corti} provides the source, data
and retained evidence. Its Zenodo series DOI identifies the release history;
individual deposits identify fixed versions.

\section{Conclusion}
The locally validated prefix of A358321 now contains 17 terms, indexed
$3$ through $19$, with four terms beyond the source baseline.
The corresponding extension of A358862 is an exact consequence of monotonicity.
The last minimum illustrates how an elementary valuation obstruction can
complete a finite computational proof by excluding entire unfinished
prime-power subtrees. Its negative part still relies on the recorded
producer coverage of the other subtrees. The empty final interval is a
separate finite statement for order $19$; the next undetermined term
in this work is $a(20)$.

\section*{Data availability}
The source, build specifications, complete term table, coverage records,
positive certificates, arithmetic exclusions, validation programs, original
inputs and execution outputs are available through the GitHub repository and
Zenodo release history cited in References. This paper is supplied as a
monolithic TeX source and matching PDF with the repository package. The
Zenodo citation uses the series (concept) DOI; individual deposits in its
version history identify fixed releases.

\section*{AI use disclosure}
Large language models were used as research and drafting assistants during
the A358321 workflow, including OpenAI Codex and external LLM-based review
systems presented as ChatGPT, Claude, DeepSeek, Gemini, Grok, Kimi, Qwen,
MiniMax and GLM/Z.ai. These are recorded service or family attributions;
exact backend identities and model versions were not consistently
authenticated. One manuscript review was obtained through a Mistral-hosted
service and identified its model family as GLM; that backend attribution
was not independently authenticated. Meta AI was also used as a service in preliminary analysis;
its underlying model family is not assumed. The assistance and review
reports are retained privately.

\paragraph{Mathematical analysis.}
AI assistance was used to examine the distinct-factor predicate, source
relationships, finite enumeration bounds and arithmetic subtree exclusions.
The accepted statements are supported by the defining identities, preserved
OEIS source records, explicit proofs, result tables and Python/PARI arithmetic
certificates. Repeated source values in an AI response do not constitute
independent confirmation.

\paragraph{Algorithm and code.}
Codex assisted implementation and validation design. Reports identifying
Claude and MiniMax contributed to code review. Final computational claims
are tied to the retained source/build identities, recorded coverage, small
exact-set comparisons, cursor/recovery tests and independent positive
certificates. These checks do not turn LLM review or positive arithmetic
certification into a second exhaustive negative enumeration. The arithmetic
prefix exclusions were applied after the search and are not attributed to
an earlier implementation feature.

\paragraph{Manuscript.}
Codex assisted drafting, literature verification, arithmetic cross-checks and
internal examination. External LLM manuscript reviews of the initial draft
and focused reviews of the revised draft were assessed against the canonical
term table, coverage records, original
source exports, certificate programs and execution outputs. Their reported
arithmetic checks and static prefix-map replays are distinct from an
independent exhaustive replay of the negative search. Final numerical
claims rely on those local artifacts and validation, not on conversational
assertions. These reports do not establish independent reproduction of the full
negative enumeration or acceptance by a publication venue.

The author, Carlo Corti, takes full scientific responsibility for the
mathematical statements, algorithm, implementation, numerical results,
validation boundaries and conclusions of this paper.

\clearpage
\begin{thebibliography}{9}
\sloppy\raggedright
\bibitem{oeis321}
OEIS Foundation Inc., \emph{A358321: a(n) is the index of the smallest
n-gonal number with exactly n distinct prime factors},
The On-Line Encyclopedia of Integer Sequences.
Submitted by Ilya Gutkovskiy, 12 December 2022.
\url{https://oeis.org/A358321}.
Official export consulted 9 November 2026; sequence revision \#67 dated 16 February 2025.

\bibitem{oeis862}
OEIS Foundation Inc., \emph{A358862: a(n) is the smallest n-gonal number
with exactly n distinct prime factors},
The On-Line Encyclopedia of Integer Sequences.
Submitted by Ilya Gutkovskiy, 3 December 2022; terms $12$--$15$ and PARI
programs by Daniel Suteu, 4 December 2022; Mathematica program by
Harvey P. Dale, 9 September 2023.
\url{https://oeis.org/A358862}.
Official export consulted 9 November 2026; sequence revision \#17 dated 16 February 2025.

\bibitem{oeis014}
OEIS Foundation Inc., \emph{A359014: a(n) is the index of the smallest
n-gonal number with exactly n prime factors (counted with multiplicity)},
The On-Line Encyclopedia of Integer Sequences.
Submitted by Ilya Gutkovskiy, 12 December 2022; PARI program and terms
$29$--$34$ by Daniel Suteu, 12 December 2022.
\url{https://oeis.org/A359014}.
Official export consulted 9 November 2026; sequence revision \#7 dated 16 February 2025.

\bibitem{polygonal}
E. W. Weisstein, \emph{Polygonal Number}, MathWorld---A Wolfram Resource.
\url{https://mathworld.wolfram.com/PolygonalNumber.html}.
Consulted 9 November 2026.

\bibitem{distinct}
E. W. Weisstein, \emph{Distinct Prime Factors}, MathWorld---A Wolfram Resource.
\url{https://mathworld.wolfram.com/DistinctPrimeFactors.html}.
Consulted 9 November 2026.

\bibitem{pari}
The PARI Group, \emph{PARI/GP: Catalogue of GP/PARI Functions,
Arithmetic functions}, official online manual; computations here used
PARI/GP 2.15.4.
\url{https://pari.math.u-bordeaux.fr/dochtml/html-stable/Arithmetic_functions.html}.
Consulted 9 November 2026.

\bibitem{gcc}
Free Software Foundation, \emph{Using the GNU Compiler Collection:
128-bit Integers}; \emph{GNU libgomp: Enabling OpenMP}, online documentation.
\url{https://gcc.gnu.org/onlinedocs/gcc/_005f_005fint128.html};
\url{https://gcc.gnu.org/onlinedocs/libgomp/Enabling-OpenMP.html}.
Consulted 9 November 2026; production compiler GCC 13.3.0.

\bibitem{corti}
C. Corti, \emph{Extending A358321 by prime-power enumeration and arithmetic exclusions},
GitHub repository and Zenodo archive, 2026.
\url{https://github.com/carcorti/A358321};
\url{https://doi.org/10.5281/zenodo.24103849}.
Zenodo series (concept) DOI; individual records identify fixed releases.
\end{thebibliography}
\end{document}
